% !TEX root = ../main.tex
%% Corrigé du devoir surveillé n° 5 — vérifié par verifications/eval_ds.py.

\corrigeSujetEval{DS5}{Devoir surveillé \numero5}

\begin{corrigeEval}{6 points}{Permutations et déterminants}
\begin{questions}
  \item Orbites : $1\to3\to7\to4\to1$, $2\to5\to2$, $6\to8\to6$, d'où
    $\res{\sigma=(1\ 3\ 7\ 4)(2\ 5)(6\ 8)}$. Signature
    $(-1)^{3}(-1)(-1)=\res{-1}$ ; ordre $\operatorname{ppcm}(4,2,2)=\res{4}$.
    \pts{1,5 pt}
  \item Après $C_1\leftarrow C_1+\dots+C_n$, la première colonne vaut
    $a+(n-1)b$ partout : on la factorise. Puis $L_i\leftarrow L_i-L_1$
    ($i\geqslant2$) donne une matrice triangulaire de diagonale
    $1,a-b,\dots,a-b$ : $\res{D_n=\bigl(a+(n-1)b\bigr)(a-b)^{n-1}}$.
    \pts{2,5 pts}
  \item $a=0$, $b=1$ : $\det(J_n-I_n)=(n-1)(-1)^{n-1}\neq0$ pour
    $n\geqslant2$ : \res{\text{inversible}}. \pts{1 pt}
  \item $L_2+L_3=(a+b+c)\,L_1$ : les lignes sont liées, le déterminant est
    \res{\text{nul}}. \pts{1 pt}
\end{questions}
\end{corrigeEval}

\begin{corrigeEval}{8 points}{Intégration}
\begin{questions}
  \item $I_0=\bigl[\arctan x\bigr]_0^{1}=\res{\dfrac\pi4}$ et
    $I_1=\Bigl[\dfrac12\ln(1+x^{2})\Bigr]_0^{1}=\res{\dfrac{\ln2}{2}}$.
    \pts{1 pt}
  \item Sur $\intFF01$, $0\leqslant x^{n+1}\leqslant x^{n}$ : par
    croissance de l'intégrale, $(I_n)$ est positive et décroissante ; et
    $\dfrac{x^{n}}{1+x^{2}}\leqslant x^{n}$ donne
    $I_n\leqslant\dfrac{1}{n+1}$. \pts{1,5 pt}
  \item On simplifie par $1+x^{2}$ :
    \[I_n+I_{n+2}=\int_0^{1}\frac{x^{n}(1+x^{2})}{1+x^{2}}\,\mathrm{d}x
      =\int_0^{1}x^{n}\,\mathrm{d}x=\frac{1}{n+1}.\]
    Par décroissance, $2I_{n+2}\leqslant I_n+I_{n+2}\leqslant2I_n$ ; en
    appliquant aussi cet encadrement au rang $n-2$, on obtient
    \[\frac{1}{2(n+1)}\leqslant I_n\leqslant\frac{1}{2(n-1)}\qquad(n\geqslant2).\]
    Multipliées par $2n$, les deux bornes tendent vers $1$ :
    $\res{I_n\sim\dfrac{1}{2n}}$. \pts{2,5 pts}
  \item On reconnaît une somme de Riemann de la fonction continue
    $x\mapsto1/\sqrt{1+x}$ sur $\intFF01$ :
    \[\sum_{k=1}^{n}\frac{1}{\sqrt{n^{2}+kn}}
      =\frac1n\sum_{k=1}^{n}\frac{1}{\sqrt{1+k/n}}
      \xrightarrow[n\to+\infty]{}\int_0^{1}\frac{\mathrm{d}x}{\sqrt{1+x}}
      =\bigl[2\sqrt{1+x}\bigr]_0^{1}=\res{2(\sqrt2-1)}.\]
    \pts{1,5 pt}
  \item $f:x\mapsto\ln(1+x)$ est de classe $\mathcal{C}^{3}$ sur $\R_{+}$,
    et $f'''(x)=\dfrac{2}{(1+x)^{3}}\leqslant2$. L'inégalité de
    Taylor-Lagrange à l'ordre $2$ entre $0$ et $x$ donne
    \[\Bigl|\ln(1+x)-x+\frac{x^{2}}{2}\Bigr|\leqslant\frac{2}{3!}\,x^{3}
      =\res{\frac{x^{3}}{3}}.\]
    \pts{1,5 pt}
\end{questions}
\begin{center}\begin{minipage}{\linewidth}\centering
\begin{tikzpicture}[xscale=6,yscale=3,line cap=round]
  \foreach \k in {1,...,8}{
    \pgfmathsetmacro{\hh}{1/sqrt(1+\k/8)}
    \filldraw[fill=otGreenDark!14,draw=otGreenDark!70,line width=0.4pt]
      ({(\k-1)/8},0) rectangle ({\k/8},\hh);}
  \draw[otMuted,line width=0.5pt,-{Stealth[length=1.5mm]}] (-0.04,0) -- (1.12,0);
  \draw[otMuted,line width=0.5pt,-{Stealth[length=1.5mm]}] (0,-0.06) -- (0,1.2);
  \draw[otNavy,line width=1.1pt] plot[domain=0:1.06,samples=60] ({\x},{1/sqrt(1+\x)});
  \node[font=\scriptsize,anchor=north] at (0,0) {$0$};
  \node[font=\scriptsize,anchor=north] at (1,0) {$1$};
  \node[font=\scriptsize,anchor=east] at (0,1) {$1$};
  \draw[otMuted,line width=0.3pt] (-0.012,1) -- (0.012,1);
  \node[font=\scriptsize,text=otNavy,anchor=south west] at (0.5,0.86) {$y=\dfrac{1}{\sqrt{1+x}}$};
\end{tikzpicture}\\*[0.5mm]
\legende{La somme de la question 4 est l'aire des $n$ rectangles (ici
$n=8$) : elle tend vers l'aire sous la courbe, $2(\sqrt2-1)\approx0{,}828$.}
\end{minipage}\end{center}
\end{corrigeEval}

\begin{corrigeEval}{6 points}{Dénombrement}
\begin{questions}
  \item $11$ lettres : $4$ I, $4$ S, $2$ P, $1$ M ; on choisit les places
    des I, puis des S, puis des P (le M prend la dernière) :
    \[\binom{11}{4}\binom{7}{4}\binom{3}{2}=\frac{11!}{4!\,4!\,2!}=\res{34\,650}.\]
    \pts{1,5 pt}
  \item \emph{Full} : hauteur du brelan ($8$), ses trois couleurs
    ($\binom43$), hauteur de la paire ($7$), ses couleurs ($\binom42$) :
    $8\cdot4\cdot7\cdot6=\res{1344}$. \emph{Double paire} : les deux
    hauteurs ($\binom82$), les couleurs ($\binom42^{2}$), la cinquième
    carte parmi les $24$ cartes des six autres hauteurs : $28\cdot36\cdot24=\res{24\,192}$.
    \pts{2,5 pts}
  \item On donne d'abord un bonbon à chacun, puis on répartit librement
    les $6$ restants (étoiles et barres, chapitre 16) :
    $\binom{6+3}{3}=\res{84}$. Sans contrainte : $\binom{13}{3}=\res{286}$.
    \pts{2 pts}
\end{questions}
\end{corrigeEval}
